\section{Artifact-Grounded Peer Review}

\ra\ uses two complementary review modes. \textbf{PDF-only review} evaluates the final manuscript as a standalone paper. In our study, \sar\ provides this view: it receives the paper and returns an ICLR-style score. \textbf{Artifact-grounded review} evaluates the paper together with the code, execution logs, result artifacts, and references. This mode evaluates paper-artifact faithfulness, not just paper plausibility.

\begin{table}[t]
\centering
\small
\begin{tabularx}{\linewidth}{p{0.28\linewidth} X X}
\toprule
Evaluation target & PDF-only \sar & Artifact-grounded \pr \\
\midrule
Inputs available & Final manuscript & Manuscript plus code, configs, logs, result files, figures, and references \\
Can inspect code? & \no & \yes \\
Can inspect logs? & \no & \yes \\
Can verify result files? & \no & \yes \\
Can detect setting mismatch? & Usually no & Yes, when settings are encoded in code, configs, or logs \\
Can detect fake references? & Partly from bibliography text & Yes, with explicit reference checks \\
Main failure mode & Rewards plausible but unsupported papers & Depends on reviewer skill and can still miss subtle defects \\
\bottomrule
\end{tabularx}
\caption{Evaluation target comparison. PDF-only review and artifact-grounded peer review answer different questions. The latter audits the relationship between claims and process artifacts.}
\label{tab:target-comparison}
\end{table}

\paragraph{Rubric.}
Every paper is reviewed by all three CLI agents, producing 351 peer reviews in total. Reviewers score nine dimensions on a 0--10 ICLR-style scale: novelty, soundness, significance, clarity, reproducibility, experimental rigor, references, reference integrity, and results integrity. The overall \pr\ score is the mean of the reviewers' overall scores. We report both aggregate scores and dimension-level patterns.

\paragraph{Reviewer aggregation.}
The use of all three agents as reviewers is deliberate. Reviewer severity varies substantially: Codex is the strictest reviewer, Kimi Code is the most lenient, and Claude Code lies between them. Multi-reviewer aggregation reduces dependence on any single reviewer's calibration. It does not eliminate bias, but it exposes severity differences that would be hidden under a single-reviewer protocol.

\paragraph{Manual integrity audit.}
In addition to automatic peer review, we manually annotate every generated paper for four integrity categories: result mismatch only, setting mismatch only, both result and setting mismatch, and fake references. A result mismatch means reported numbers do not match stored outputs or logs. A setting mismatch means the paper describes a method, dataset, baseline, hyperparameter, scale, or execution setting that differs from the artifact. The ``both'' category records papers exhibiting both mismatch types. Fake references include nonexistent citations, fabricated metadata, or references that cannot be verified as stated.

\paragraph{Conceptual decomposition.}
The results suggest a useful decomposition for evaluating auto-research systems: apparent paper quality, artifact faithfulness, and process reliability. Apparent paper quality captures whether the manuscript reads like a competent submission. Artifact faithfulness captures whether the manuscript is supported by code, logs, and result files. Process reliability captures whether the agent consistently designs, executes, records, and interprets experiments without drift. A system can score well on the first component while failing the latter two; artifact-grounded review is designed to make those failures visible.
